\documentclass[12pt,a4paper]{article}
\usepackage[utf8]{inputenc}
\usepackage[russian]{babel}
\usepackage{geometry}
\geometry{top=2cm,bottom=2cm,left=2.5cm,right=2cm}
\usepackage{graphicx} 
\usepackage{amsmath}  
\usepackage{booktabs} 
\usepackage{longtable} 
\usepackage{float}
\usepackage{tikz}     
\usetikzlibrary{shapes,arrows}

\title{«Яндекс.Локер» (городские камеры хранения)}
\author{Кузовкова Милана}
\date{\today}

\begin{document}

\maketitle
\tableofcontents
\newpage

\section{Введение и постановка задачи}
\subsection{Концепция проекта}
Яндекс.Локер — это сеть автоматизированных камер хранения, расположенных в точках повышенного пешеходного трафика (метро, ТЦ, вокзалы, аэропорты). 

\textbf{Цель проекта:} оценить экономическую целесообразность запуска пилотного кластера из 10 стоек в городе-миллионнике.

\subsection{Гипотезы}
\begin{itemize}
    \item Спрос имеет ярко выраженную сезонность и привязку к погоде/праздникам.
    \item Основной драйвер выручки — не аренда, а штрафы за просрочку (как в каршеринге).
    \item Известный бренд «Яндекс» снижает CAC на 40\% относительно конкурентов.
\end{itemize}

\subsection{Методология}
Расчет построен на:
\begin{enumerate}
    \item Юнит-экономике — расчет прибыли с 1 стойки.
    \item Моделировании загрузки.
    \item Сравнении ROI с альтернативными инвестициями.
\end{enumerate}

\section{Анализ рынка}
\begin{table}[H]
\centering
\caption{Оценка емкости рынка для города-миллионника}
\begin{tabular}{@{}lcc@{}}
\toprule
Показатель & Количество & Примечание \\
\midrule
Население города & 1 200 000 чел & \\
Туристический поток в год & 3 500 000 чел & Включая командировки \\
Потенциальные пользователи (TAM) & 650 000 чел & 40\% горожан + туристы \\
Достижимый рынок (SAM) & 65 000 чел & 10\% от TAM (только ЦА) \\
Реальный план (SOM) & 6 500 чел & 10\% от SAM (пилот) \\
\bottomrule
\end{tabular}
\end{table}

\textit{Вывод:} даже при захвате 1\% рынка пилотный проект обслуживает 6500 уникальных клиентов в год.

\footnotetext{%
    TAM (Total Addressable Market) - общий объём рынка, который теоретически может быть охвачен продуктом. Включает всех потенциальных потребителей без учёта ограничений. 
    
    SAM (Serviceable Available Market) - часть TAM, которую компания может реально обслужить с учётом географических, технических и ценовых ограничений. 
    
    SOM (Serviceable Obtainable Market) - часть SAM, которую компания планирует захватить с учётом конкуренции и маркетинговых возможностей.}

\section{Инвестиционные затраты}
\subsection{Оборудование}
\begin{itemize}
    \item Умная стойка (8 ячеек, 3 размера) $\sim$ 850 000 руб.
    \item Доставка и монтаж (10\% от стоимости): 85 000 руб.
    \item Интеграция с приложением API Яндекс $\sim$ в 300 000 руб (разовые работы).
\end{itemize}

\subsection{Итоговый CAPEX на 10 стоек}
\begin{equation}
CAPEX = 10 \times (850\,000 + 85\,000) + 300\,000 = 9\,650\,000 \text{ руб.}
\end{equation}

\section{Операционные затраты в месяц}
\begin{table}[H]
\centering
\caption{Ежемесячные операционные расходы}
\begin{tabular}{@{}lcc@{}}
\toprule
Показатель & На 1 стойку & На 10 стоек \\
\midrule
Аренда места & 25 000 & 250 000 \\
Электроэнергия и интернет & 3 000 & 30 000 \\
Техобслуживание (удаленно) & 5 000 & 50 000 \\
Страховка вещей & 2 000 & 20 000 \\
Комиссия эквайринга (2\%) & 605 & 6 050 \\
ИТ-поддержка  & - & 40 000 \\
\bottomrule
\end{tabular}
\end{table}

Комиссия эквайринга составляет 2\% от выручки, включая штрафы за просрочку.

Выручка с одной стойки при загрузке 60\% с учётом штрафов:
\[
25\,200 \times 1.2 = 30\,240 \text{ руб./мес.}
\]

Комиссия эквайринга с одной стойки:
\[
30\,240 \times 0.02 = 604.8 \approx 605 \text{ руб./мес.}
\]

Для 10 стоек:
\[
605 \times 10 = 6\,050 \text{ руб./мес.}
\]

\textit{Примечание:} комиссия рассчитана исходя из предположения, что все платежи (включая штрафы) проходят через безналичную оплату.

\[
OPEX_{фикс} = 250\,000 + 30\,000 + 50\,000 + 20\,000 + 6\,050 + 40\,000 = 396\,050 \text{ руб./мес.}
\]

\[
OPEX_{на\,стойку} = \frac{396\,050}{10} = 39\,605 \text{ руб./мес.}
\]

\footnotetext{CAPEX (Capital Expenditure) - капитальные затраты, единовременные расходы на приобретение или создание внеоборотных активов (оборудование, монтаж, разработка ПО). \\
OPEX (Operational Expenditure) - операционные затраты, регулярные расходы, необходимые для поддержания деятельности (аренда, электроэнергия, обслуживание, зарплата). \\
API (Application Programming Interface) - интерфейс программирования приложений, набор правил и инструментов для интеграции оборудования с программным обеспечением (в данном случае с приложением Яндекс).}

\section{Модель доходов}
\subsection{Структура тарифа}
\begin{itemize}
    \item Малая ячейка (рюкзак, сумка) — 150 руб./сутки.
    \item Средняя ячейка (малый чемодан, дорожная сумка) — 250 руб./сутки.
    \item Большая ячейка (большой чемодан, габаритные вещи) — 400 руб./сутки.
    \item \textbf{Штраф за просрочку} — 100 руб./час (после 24 часов).
\end{itemize}

\subsection{Расчет средней выручки с 1 стойки в день}
Допустим, загрузка ячеек:
\begin{itemize}
    \item 40\% — малые, 30\% — средние, 10\% — большие, 20\% — пустуют.
    \item Средняя заполненность: \textbf{60\%}.
\end{itemize}

\begin{equation}
ARPU_{день} = 0.4 \cdot 150 + 0.3 \cdot 250 + 0.1 \cdot 400 = 60 + 75 + 40 = 175 \text{ руб. с одной ячейки}
\end{equation}

Для стойки из 8 ячеек:
\[
Revenue_{день} = 8 \times 175 \times 0.6 = 840 \text{ руб./день}
\]
\[
Revenue_{мес} = 840 \times 30 = 25\,200 \text{ руб./мес. с одной стойки}
\]

\subsection{Учет штрафов }
По статистике каршеринга, около 20\% клиентов задерживаются. Добавим +20\% к выручке.
\[
Revenue_{мес}^{скорр} = 25\,200 \times 1.2 = 30\,240 \text{ руб.}
\]

\section{Расчет окупаемости}
\subsection{Прибыль с одной стойки}
\[
EBITDA_{мес} = 30\,240 - (396\,050 / 10) = 30\,240 - 39\,605 = -9\,365 \text{ руб. (убыток!)}
\]

\textbf{Проблема:} при загрузке 60\% стойка убыточна. Найдем точку безубыточности.

\subsection{Точка безубыточности}
\[
\text{Нужная выручка} = OPEX_{на\,стойку} = 39\,605 \text{ руб./мес.}
\]
\[
\text{Необходимая загрузка} = \frac{39\,605}{8 \times 175 \times 1.2 \times 30} \approx 79\%
\]

\textit{Вывод:} каждая стойка должна быть заполнена минимум на 79\%, чтобы выйти в ноль.

\subsection{Окупаемость инвестиций}
\[
Payback = \frac{CAPEX}{EBITDA_{мес} \times 12} = \frac{965\,000}{(-9\,365) \times 12} = \text{бесконечность (убыток)}
\]

Следовательно, бизнес нежизнеспособен при текущих ценах.

Необходимо:
\begin{enumerate}
    \item Повысить тариф на 30\%.
    \item Снизить аренду за счет переговоров с ТЦ (revenue share вместо фиксированной).
    \item Увеличить долю штрафов (маркетинг на забывчивость).
\end{enumerate}

\section{Визуализация }
\begin{figure}[H]
\centering
\begin{tikzpicture}
    \draw[->] (0,0) -- (6,0) node[right] {Загрузка, \%};
    \draw[->] (0,0) -- (0,5) node[above] {Прибыль, тыс. руб.};
    \draw[red, thick] (0,2.5) -- (6,2.5) node[right] {Точка безубыточности};
    \draw[blue, thick] (0,1) -- (3,1.5) -- (5,3) node[right] {Фактическая прибыль};
\end{tikzpicture}
\caption{График зависимости прибыли от загрузки}
\end{figure}

\section{Заключение }
\begin{enumerate}
    \item Пилотный проект не окупается при стандартных тарифах. Рекомендуемый тариф от 350 руб./сутки.
    \item Самый рискованный показатель - аренда места. Нужно договариваться о проценте от выручки (10-15\%).
    \item Проект имеет смысл только при загрузке > 80\% и в местах с туристическим трафиком (вокзалы).
\end{enumerate}


\newpage
\section*{Словарь терминов}
\addcontentsline{toc}{section}{Словарь терминов}

\begin{longtable}{@{}p{0.25\textwidth}p{0.70\textwidth}@{}}
\toprule
\textbf{Термин} & \textbf{Определение} \\
\midrule
\endfirsthead
\toprule
\textbf{Термин} & \textbf{Определение} \\
\midrule
\endhead
\bottomrule
\endfoot

\textbf{API} & 
Application Programming Interface - интерфейс программирования приложений, набор правил и инструментов для интеграции оборудования с программным обеспечением. \\
\midrule

\textbf{ARPU} & 
Average Revenue Per User - средняя выручка на одного пользователя (в данном проекте — на одну ячейку в день). \\
\midrule

\textbf{CAC} & 
Customer Acquisition Cost - стоимость привлечения одного клиента. \\
\midrule

\textbf{CAPEX} & 
Capital Expenditure - капитальные затраты, единовременные расходы на приобретение или создание внеоборотных активов (оборудование, монтаж, разработка ПО). \\
\midrule

\textbf{EBITDA} & 
Earnings Before Interest, Taxes, Depreciation and Amortization - прибыль до вычета процентов, налогов и амортизации. \\
\midrule

\textbf{OPEX} & 
Operational Expenditure - операционные затраты, регулярные расходы, необходимые для поддержания деятельности (аренда, электроэнергия, обслуживание, зарплата). \\
\midrule

\textbf{Revenue} & 
Выручка — общая сумма денег, полученная от реализации услуг. \\
\midrule

\textbf{ROI} & 
Return on Investment - коэффициент возврата инвестиций. \\
\midrule

\textbf{SAM} & 
Serviceable Available Market - часть TAM, которую компания может реально обслужить с учётом географических, технических и ценовых ограничений. \\
\midrule

\textbf{SOM} & 
Serviceable Obtainable Market - часть SAM, которую компания планирует захватить с учётом конкуренции и маркетинговых возможностей. \\
\midrule

\textbf{TAM} & 
Total Addressable Market - общий объём рынка, который теоретически может быть охвачен продуктом.\\
\midrule

\end{longtable}
\end{document}